An exact algebraic number is represented by a square-free integer polynomial, a rational isolating interval containing exactly one real root, and a Thom-sign record.  Exact field operations are represented by resultant identities followed by isolating-interval refinement.  A sign certificate for a polynomial at algebraic arguments contains the relevant subresultant or Sturm remainder sequence and rational endpoint signs; zero is certified by an exact remainder or greatest-common-divisor identity.  Every primitive checker uses only integer and rational arithmetic.

An \(\mathsf{AlgTrace}\) contains: a finite input formula over an effectively presented real-algebraic field; the output of the fixed quantifier-elimination algorithm; the projection-polynomial lists of the fixed adapted cylindrical-decomposition algorithm; for each lifting stack, complete root-isolation traces and section or sector sample records; exact sign tables; recursive finite truth tables for quantified blocks; and, for optimization, the one-dimensional image formula, its ordered true cells, algebraic endpoint records, and separate nonempty-fiber lifting traces for every reported endpoint.  The checker \(\operatorname{CheckAlg}\) recomputes polynomial normalization, resultants, Sturm variations, root counts, signs, finite disjunctions for existential blocks, finite conjunctions for universal blocks, endpoint order, feasibility substitutions, and objective substitutions.  It never accepts a bare assertion of formula equivalence, nonemptiness, optimality, or witness existence.

The hybrid terminal grammar is
\[
\begin{split}
\mathsf{HybridResult}::={}&
\operatorname{ID\_RAT}(F,\delta)
\mid\operatorname{ID\_CAN}(\alpha,\tau)
\mid\operatorname{PAIR\_CAN}(\vartheta_-,\vartheta_+,\tau)\\
&\mid\operatorname{ZERO\_EVIDENCE}(\vartheta_0,\tau)
\mid\operatorname{NOT\_UNIFORMLY\_DEFINED}(\vartheta_0,\vartheta_+,\tau).
\end{split}
\]
The deterministic algorithm \(\mathsf{DCTFTR\text{-}CAN}\) first runs \(\mathsf{DCTFTR}\).  It relabels a successful rational identification or uniform-zero trace.  Every other terminal invokes the fixed canonical polynomial builder on \(\Theta_G^{\mathrm{can}}(\mathcal P)\), the fixed quantifier-elimination and adapted-decomposition algorithms, and the exact checker above.  It first computes the evidence range; it rejects division unless the checked lower endpoint is strictly positive, and only then computes the conditional-ratio range.  Canonical finite orders and fixed algorithms make the selected run unique, but no uniqueness of CAD representation or certificate is asserted.